\documentclass[10pt,a4paper]{amsart}
\usepackage{amsmath,amssymb,amsthm,graphicx}
\usepackage{verbatim}
\newenvironment{reference}{\comment}{\endcomment}
%\newenvironment{reference}{}{}
\newenvironment{slogan}{\comment}{\endcomment}
\newenvironment{history}{\comment}{\endcomment}

% For commutative diagrams we use Xy-pic
\usepackage[all]{xy}

% We use 2cell for 2-commutative diagrams.
\xyoption{2cell}
% Theorem environments.
%
\theoremstyle{plain}
\newtheorem{theorem}[subsection]{Theorem}
\newtheorem{proposition}[subsection]{Proposition}
\newtheorem{lemma}[subsection]{Lemma}

\theoremstyle{definition}
\newtheorem{definition}[subsection]{Definition}
\newtheorem{example}[subsection]{Example}
\newtheorem{exercise}[subsection]{Exercise}
\newtheorem{situation}[subsection]{Situation}

\theoremstyle{remark}
\newtheorem{remark}[subsection]{Remark}
\newtheorem{remarks}[subsection]{Remarks}

\numberwithin{equation}{subsection}

% Macros
%
\def\lim{\mathop{\mathrm{lim}}\nolimits}
\def\colim{\mathop{\mathrm{colim}}\nolimits}
\def\Spec{\mathop{\mathrm{Spec}}}
\def\Hom{\mathop{\mathrm{Hom}}\nolimits}
\def\Ext{\mathop{\mathrm{Ext}}\nolimits}
\def\SheafHom{\mathop{\mathcal{H}\!\mathit{om}}\nolimits}
\def\SheafExt{\mathop{\mathcal{E}\!\mathit{xt}}\nolimits}
\def\Sch{\mathit{Sch}}
\def\Mor{\mathop{\mathrm{Mor}}\nolimits}
\def\Ob{\mathop{\mathrm{Ob}}\nolimits}
\def\Sh{\mathop{\mathit{Sh}}\nolimits}
\def\NL{\mathop{N\!L}\nolimits}
\def\CH{\mathop{\mathrm{CH}}\nolimits}
\def\proetale{{pro\text{-}\acute{e}tale}}
\def\etale{{\acute{e}tale}}
\def\QCoh{\mathit{QCoh}}
\def\Ker{\mathop{\mathrm{Ker}}}
\def\Im{\mathop{\mathrm{Im}}}
\def\Coker{\mathop{\mathrm{Coker}}}
\def\Coim{\mathop{\mathrm{Coim}}}

% Boxtimes
%
\DeclareMathSymbol{\boxtimes}{\mathbin}{AMSa}{"02}

%
% Macros for moduli stacks/spaces
%
\def\QCohstack{\mathcal{QC}\!\mathit{oh}}
\def\Cohstack{\mathcal{C}\!\mathit{oh}}
\def\Spacesstack{\mathcal{S}\!\mathit{paces}}
\def\Quotfunctor{\mathrm{Quot}}
\def\Hilbfunctor{\mathrm{Hilb}}
\def\Curvesstack{\mathcal{C}\!\mathit{urves}}
\def\Polarizedstack{\mathcal{P}\!\mathit{olarized}}
\def\Complexesstack{\mathcal{C}\!\mathit{omplexes}}
% \Pic is the operator that assigns to X its picard group, usage \Pic(X)
% \Picardstack_{X/B} denotes the Picard stack of X over B
% \Picardfunctor_{X/B} denotes the Picard functor of X over B
\def\Pic{\mathop{\mathrm{Pic}}\nolimits}
\def\Picardstack{\mathcal{P}\!\mathit{ic}}
\def\Picardfunctor{\mathrm{Pic}}
\def\Deformationcategory{\mathcal{D}\!\mathit{ef}}
\usepackage{cleveref}
\crefname{originaldefinition}{original definition}{original definitions}
\Crefname{originaldefinition}{Original definition}{Original definitions}
\makeatletter
\newlabel{divisors-definition-regular-ideal-sheaf}{{063D}{source}{}{sourcepointer-1}{}}
\newlabel{divisors-definition-regular-ideal-sheaf@cref}{{[originaldefinition][0][]063D}{[1][0][]source}{}{}{}}
\makeatother
\crefname{originallemma}{original lemma}{original lemmas}
\Crefname{originallemma}{Original lemma}{Original lemmas}
\makeatletter
\newlabel{cohomology-lemma-internal-hom-composition}{{0A8V}{source}{}{sourcepointer-2}{}}
\newlabel{cohomology-lemma-internal-hom-composition@cref}{{[originallemma][0][]0A8V}{[1][0][]source}{}{}{}}
\makeatother
\crefname{originalequation}{original equation}{original equations}
\Crefname{originalequation}{Original equation}{Original equations}
\makeatletter
\newlabel{more-algebra-equation-simple-tor-product}{{068H}{source}{}{sourcepointer-3}{}}
\newlabel{more-algebra-equation-simple-tor-product@cref}{{[originalequation][0][]068H}{[1][0][]source}{}{}{}}
\makeatother
\crefname{originallemma}{original lemma}{original lemmas}
\Crefname{originallemma}{Original lemma}{Original lemmas}
\makeatletter
\newlabel{cohomology-lemma-Rhom-strictly-perfect}{{08DM}{source}{}{sourcepointer-4}{}}
\newlabel{cohomology-lemma-Rhom-strictly-perfect@cref}{{[originallemma][0][]08DM}{[1][0][]source}{}{}{}}
\makeatother
\title{For a Koszul regular ideal the Ext algebra of the quotient is the exterior algebra on its first Ext sheaf}
\author{The Stacks Project Authors}
\date{stacks-project:duality.tex:local-snapshot-20261012; GNU FDL 1.2 or later}
\begin{document}
\maketitle
\noindent\textit{Editorial note.} All delivered mathematical text is the author's own, copied from the pinned Stacks Project chapter file duality.tex; no statement or proof step was rewritten, reordered or removed.\par\medskip
\noindent\textit{Editorial note.} Preamble provenance: the pinned chapter file opens with an input of the shared preamble and carries no preamble text of its own, so the project's single Stacks-wide preamble staged at sources/stacks-duality/preamble.tex is excerpted verbatim under the recipe fixed by LANE\_BRIEF 22.1, as three spans, lines 8-13 (verbatim package and the reference/slogan/history environments), lines 14-18 (Xy-pic setup) and lines 154-217 (the native theorem declarations and mathematical macros). Nothing was added, removed or redefined; the wrapper class and the amsmath/amssymb/amsthm/graphicx base packages are supplied by the assembler.\par\medskip
\noindent\textit{Editorial note.} The counted statement and proof contain no cross reference, no citation and no included file, so the delivered body is self contained.\par\medskip
\noindent\textit{Editorial note.} Canonical location of the counted proposition: Stacks tag 0BQY (sources/stacks-duality/pinned-original-source.tex line 3661).\par\medskip
\noindent\textit{Allowed standard prerequisites (editorial).}\begin{itemize}
\item Koszul regular sheaves of ideals. Divisors, Definition labelled divisors-definition-regular-ideal-sheaf, invoked in the statement. Supplies the hypothesis and the local Koszul regular sequence.
\item Composition of internal Homs and the simple tensor product comparison. Cohomology, Lemma labelled cohomology-lemma-internal-hom-composition, cited at proof lines 3678-3680, and More on Algebra, Equation labelled more-algebra-equation-simple-tor-product, cited at proof lines 3688-3692. Fixes the multiplication whose exterior-algebra structure is computed.
\item Derived Hom of a strictly perfect complex. Cohomology, Lemma labelled cohomology-lemma-Rhom-strictly-perfect, cited at proof lines 3714-3718. Replaces RSheafHom by the Hom complex of the Koszul dga with vanishing differentials.
\item The Koszul complex as a resolution of the quotient. The quasi-isomorphism A\textasciicircum{}bullet -> O\_X built from the divided power polynomial algebra in proof lines 3700-3712. Identifies the Ext sheaves with the Hom sheaves out of the graded pieces A\textasciicircum{}\{-i\}.
\end{itemize}\medskip
\ifdefined\hypertarget\hypertarget{sourcepointer-1}{}\fi
\noindent\textit{Source reference.} Original definition 063D: sources/stacks-chow/tags.txt, source lines 7827--7827. Cross-chapter standard prerequisite located in the project tag list staged in the local snapshot; invoked at proof line None of the counted proof. It is retained only as a navigation pointer; every load-bearing step is inside the delivered span.\par\medskip
\ifdefined\hypertarget\hypertarget{sourcepointer-2}{}\fi
\noindent\textit{Source reference.} Original lemma 0A8V: sources/stacks-chow/tags.txt, source lines 12928--12928. Cross-chapter standard prerequisite located in the project tag list staged in the local snapshot; invoked at proof line 3685 of the counted proof. It is retained only as a navigation pointer; every load-bearing step is inside the delivered span.\par\medskip
\ifdefined\hypertarget\hypertarget{sourcepointer-3}{}\fi
\noindent\textit{Source reference.} Original equation 068H: sources/stacks-chow/tags.txt, source lines 8006--8006. Cross-chapter standard prerequisite located in the project tag list staged in the local snapshot; invoked at proof line 3695 of the counted proof. It is retained only as a navigation pointer; every load-bearing step is inside the delivered span.\par\medskip
\ifdefined\hypertarget\hypertarget{sourcepointer-4}{}\fi
\noindent\textit{Source reference.} Original lemma 08DM: sources/stacks-chow/tags.txt, source lines 10642--10642. Cross-chapter standard prerequisite located in the project tag list staged in the local snapshot; invoked at proof line 3726 of the counted proof. It is retained only as a navigation pointer; every load-bearing step is inside the delivered span.\par\medskip
\begin{lemma}
\label{lemma-regular-ideal-ext}
Let $Y$ be a ringed space. Let $\mathcal{I} \subset \mathcal{O}_Y$
be a sheaf of ideals. Set $\mathcal{O}_X = \mathcal{O}_Y/\mathcal{I}$.
If $\mathcal{I}$ is Koszul-regular
(Divisors, Definition \ref{divisors-definition-regular-ideal-sheaf})
then composition on $R\SheafHom_{\mathcal{O}_Y}(\mathcal{O}_X, \mathcal{O}_X)$
defines isomorphisms
$$
\wedge^i(\SheafExt^1_{\mathcal{O}_Y}(\mathcal{O}_X, \mathcal{O}_X))
\longrightarrow
\SheafExt^i_{\mathcal{O}_Y}(\mathcal{O}_X, \mathcal{O}_X)
$$
for all $i$.
\end{lemma}

\begin{proof}
By composition we mean the map
$$
R\SheafHom_{\mathcal{O}_Y}(\mathcal{O}_X, \mathcal{O}_X)
\otimes_{\mathcal{O}_Y}^\mathbf{L}
R\SheafHom_{\mathcal{O}_Y}(\mathcal{O}_X, \mathcal{O}_X)
\longrightarrow
R\SheafHom_{\mathcal{O}_Y}(\mathcal{O}_X, \mathcal{O}_X)
$$
of Cohomology, Lemma \ref{cohomology-lemma-internal-hom-composition}.
This induces multiplication maps
$$
\SheafExt^a_{\mathcal{O}_Y}(\mathcal{O}_X, \mathcal{O}_X)
\otimes_{\mathcal{O}_Y}
\SheafExt^b_{\mathcal{O}_Y}(\mathcal{O}_X, \mathcal{O}_X)
\longrightarrow
\SheafExt^{a + b}_{\mathcal{O}_Y}(\mathcal{O}_X, \mathcal{O}_X)
$$
Please compare with
More on Algebra, Equation (\ref{more-algebra-equation-simple-tor-product}).
The statement of the lemma means that the induced map
$$
\SheafExt^1_{\mathcal{O}_Y}(\mathcal{O}_X, \mathcal{O}_X)
\otimes \ldots \otimes
\SheafExt^1_{\mathcal{O}_Y}(\mathcal{O}_X, \mathcal{O}_X)
\longrightarrow
\SheafExt^i_{\mathcal{O}_Y}(\mathcal{O}_X, \mathcal{O}_X)
$$
factors through the wedge product and then induces an isomorphism.
To see this is true we may work locally on $Y$. Hence we may assume
that we have global sections $f_1, \ldots, f_r$ of $\mathcal{O}_Y$
which generate $\mathcal{I}$ and which form a Koszul regular sequence.
Denote
$$
\mathcal{A} = \mathcal{O}_Y\langle \xi_1, \ldots, \xi_r\rangle
$$
the sheaf of strictly commutative differential graded $\mathcal{O}_Y$-algebras
which is a (divided power) polynomial algebra on
$\xi_1, \ldots, \xi_r$ in degree $-1$ over $\mathcal{O}_Y$
with differential $\text{d}$ given by the rule $\text{d}\xi_i = f_i$.
Let us denote $\mathcal{A}^\bullet$ the underlying
complex of $\mathcal{O}_Y$-modules which is the Koszul complex
mentioned above. Thus the canonical map
$\mathcal{A}^\bullet \to \mathcal{O}_X$
is a quasi-isomorphism. We obtain quasi-isomorphisms
$$
R\SheafHom_{\mathcal{O}_Y}(\mathcal{O}_X, \mathcal{O}_X) \to
\SheafHom^\bullet(\mathcal{A}^\bullet, \mathcal{A}^\bullet) \to
\SheafHom^\bullet(\mathcal{A}^\bullet, \mathcal{O}_X)
$$
by Cohomology, Lemma \ref{cohomology-lemma-Rhom-strictly-perfect}.
The differentials of the latter complex are zero, and hence
$$
\SheafExt^i_{\mathcal{O}_Y}(\mathcal{O}_X, \mathcal{O}_X)
\cong \SheafHom_{\mathcal{O}_Y}(\mathcal{A}^{-i}, \mathcal{O}_X)
$$
For $j \in \{1, \ldots, r\}$ let $\delta_j : \mathcal{A} \to \mathcal{A}$
be the derivation of degree $1$ with $\delta_j(\xi_i) = \delta_{ij}$
(Kronecker delta). A computation shows that
$\delta_j \circ \text{d} = - \text{d} \circ \delta_j$ which shows that
we get a morphism of complexes
$$
\delta_j : \mathcal{A}^\bullet \to \mathcal{A}^\bullet[1].
$$
Whence $\delta_j$ defines a section of the corresponding
$\SheafExt$-sheaf.
Another computation shows that $\delta_1, \ldots, \delta_r$
map to a basis for $\SheafHom_{\mathcal{O}_Y}(\mathcal{A}^{-1}, \mathcal{O}_X)$
over $\mathcal{O}_X$.
Since it is clear that $\delta_j \circ \delta_j = 0$
and $\delta_j \circ \delta_{j'} = - \delta_{j'} \circ \delta_j$
as endomorphisms of $\mathcal{A}$ and hence in the
$\SheafExt$-sheaves
we obtain the statement that our map above factors through
the exterior power. To see we get the desired isomorphism
the reader checks that the elements
$$
\delta_{j_1} \circ \ldots \circ \delta_{j_i}
$$
for $j_1 < \ldots < j_i$ map to a basis of the sheaf
$\SheafHom_{\mathcal{O}_Y}(\mathcal{A}^{-i}, \mathcal{O}_X)$
over $\mathcal{O}_X$.
\end{proof}

\end{document}
